\addcontentsline{toc}{chapter}{Abstract}
\chapter*{Abstract}

\noindent
\begin{center}
	\begin{minipage}{0.9\linewidth}
		\centering
		\textbf{\MakeUppercase{\titleEnglish}}
	\end{minipage}
\end{center}

\vspace{2em}

\noindent
\begin{center}
	\renewcommand{\arraystretch}{1.25}
	\begin{tabularx}{0.75\linewidth}{@{}l@{\ }l@{\enspace}X}
		Student Name / NRP &:& \theAuthor\ / \theAuthorId \\
		Departement &:& \departmentNameHeadlessEnglish\ \facultyNameEnglishAbbreviated\ - \campusNameAbbreviated \\
		Supervisor / NIP &:& \supervisor\ / \supervisorId
	\end{tabularx}
\end{center}

\vspace{1em}

\section*{Abstract}

The metric dimension and the edge metric dimension of a graph are fundamental invariants in graph theory. Roughly speaking, the metric dimension is the minimum number of vertices needed to distinguish all other vertices in the graph, while the edge metric dimension does the same for the edges. These invariants have a range of applications, such as in network navigation, source tracking, robotics, and biological systems.

This thesis explores the use of Satisfiability Modulo Theories (SMT) as an approach to compute the metric dimension and the edge metric dimension. The problems are formulated declaratively in SMT and solved using a solver. Canonical graph representations are employed to remove redundancy caused by automorphisms, and experiments are conducted on all unique connected graphs with 2 to 9 vertices, as well as on some larger graphs.

The experimental results show that all SMT formulations produce correct values for metric dimension and edge metric dimension. In terms of execution time, for small-sized graphs, the brute-force approach is faster than SMT. In contrast, for large-sized graphs, the SMT approach with pseudo boolean theory exhibits the best execution performance, while the brute-force method fails entirely to complete the computation within the given time limit. These results indicate that SMT can be used to generate correct graph invariants for metric dimension, with a practical execution-time growth as the graph size grows.

\vfill

\noindent
\begin{minipage}[t][5cm]{\textwidth}
	\begin{tabularx}{\linewidth}{@{}l@{\enspace}X}
		\textbf{Keywords:} & \textit{Metric Dimension, Edge Metric Dimension, Satisfiability Modulo Theories, SMT Solver, NP-hard}
	\end{tabularx}
\end{minipage}

%
